\begin{abstract}
Organizations that keep telemetry in-house must triage alerts with small local LLMs, and on their own these models cannot do it. Many alerts, however, belong to recurring types for which analysts already know the candidate causes, the checks that tell them apart, and how past cases were resolved. We present \hesp{}, which makes a small model usable for such alert types by dividing triage into two jobs: the model \emph{reads} each probe's raw log text into one of the outcomes analysts defined, and a controller \emph{decides} which probe to run, when to stop, and which verdict to accept, using a Bayesian ledger over the candidate causes with likelihoods counted from resolved cases. In three simulated triage settings with five open-weight models from 7B to 72B, small models deciding on their own resolve at most a third of cases even when every observation is already labeled, and none of 48 when they must read raw logs. The controller alone resolves every case in two settings when observations are structured, but when a logging change alters the log format, its rule parser reads none of 364 observations and every case escalates. Together, with Qwen2.5-7B reading and the controller deciding, all 48 drifted-format episodes are resolved, and 34 with Llama-3.1-8B. Reading is also an attack channel: an attacker who writes a consistent benign story into every log makes Qwen2.5-7B as reader close all 36 attacks as benign. Never letting model-parsed evidence support a benign verdict stops every attack we tested, at the cost of escalating honest benign cases when no trusted parser can read the logs. Our settings are simulated, their causes and outcomes are written in advance, and each cause leaves a distinctive observation; new alert types and tables counted from real tickets are outside their scope.
\end{abstract}
